%----------------------------------------------------------------------------------------
% METADATA BLOCK (YAML)  — see inbox/README.md
%----------------------------------------------------------------------------------------
% uid: ""
% school: "lycee"
% level: "2eme"
% section: ["Sciences", "Informatique"]
% chapter: "Notion de Polynômes"
% title: "Une racine double"
% difficulty: "Facile"
% tags: ["polynomes", "racine-double", "tableau-de-signe", "inequation"]
% provenance: "own"
% batch: "2026-11-19-2sc-polynomes-barycentre"
%----------------------------------------------------------------------------------------

\newcommand{\R}{\mathbb{R}}
\newcommand{\vect}[1]{\overrightarrow{#1}}

% ---STATEMENT---
Soit $P(x)=x^3-3x-2$.
\begin{enumerate}
  \item Calculer $P(-1)$ et $P(2)$.
  \item Déterminer les réels $a$, $b$, $c$ tels que, pour tout réel $x$, $P(x)=(x+1)(ax^2+bx+c)$.
  \item Montrer que $P(x)=(x+1)^2(x-2)$.
  \item Dresser le tableau de signe de $P(x)$, puis résoudre dans $\R$ l'inéquation $P(x)\geq0$.
  \item En déduire que, pour tout réel $x\leq2$ : $x^3\leq3x+2$.
\end{enumerate}

% ---SOLUTION---
\begin{enumerate}[itemsep=14pt]
  \item $P(-1)=-1+3-2=0$ \quad et \quad $P(2)=8-6-2=0$. Les réels $-1$ et $2$ sont des racines de $P$.

  \item On développe : $(x+1)(ax^2+bx+c)=ax^3+(a+b)x^2+(b+c)x+c$.

  On identifie avec $P(x)=x^3+0x^2-3x-2$ :
  \[
  \left\{
  \begin{aligned}
    a&=1 &&\text{(coefficient de }x^3\text{)}\\
    a+b&=0 &&\text{(coefficient de }x^2\text{)}\\
    b+c&=-3 &&\text{(coefficient de }x\text{)}\\
    c&=-2 &&\text{(terme constant)}
  \end{aligned}
  \right.
  \qquad\text{donc}\qquad a=1,\ b=-1,\ c=-2
  \]
  Contrôle : $b+c=-1-2=-3$ (vérifié). \quad Ainsi $P(x)=(x+1)(x^2-x-2)$.

  \item $2$ est racine de $P$ mais pas de $x+1$, donc $2$ est racine de $x^2-x-2$ : $2^2-2-2=0$.

  Le produit des racines de $x^2-x-2$ vaut $\dfrac ca=-2$, donc l'autre racine est $-1$ :
  \[x^2-x-2=(x-2)(x+1)\qquad\text{et}\qquad P(x)=(x+1)(x+1)(x-2)=(x+1)^2(x-2)\]

  \item
  \begin{prop}[Signe d'un carré]
  \ib{$(x+1)^2\geq0$} pour tout réel $x$, et $(x+1)^2=0$ seulement pour $x=-1$ : un carré ne change pas le signe d'un produit.
  \end{prop}
  \begin{center}
  
  \begin{tabular}{c|ccccccc}
    $x$ & $-\infty$ & & $-1$ & & $2$ & & $+\infty$ \\ \hline
    $(x+1)^2$ & & $+$ & $0$ & $+$ & $+$ & $+$ & \\
    $x-2$ & & $-$ & $-$ & $-$ & $0$ & $+$ & \\ \hline
    $P(x)$ & & $-$ & $0$ & $-$ & $0$ & $+$ &
  \end{tabular}
  \end{center}
  On garde les zones $+$ et les zéros :
  \[\mathcal S=\{-1\}\cup[\,2\,;+\infty[\]
  \textbf{Piège :} oublier la valeur isolée $-1$, où $P$ s'annule sans changer de signe.

  \item D'après le tableau, $P(x)\leq0$ pour tout $x\leq2$ (avec $P(-1)=0$). Donc, pour tout réel $x\leq2$ :
  \[x^3-3x-2\leq0\qquad\text{soit}\qquad x^3\leq3x+2\]
\end{enumerate}
